\begin{table}[H]
\centering
\caption{Estadísticas Descriptivas por Condición de Pobreza}
\label{tab:descriptivas}
\footnotesize % Reduce el tamaño de fuente
\begin{tabularx}{\textwidth}{@{}lXXX@{}}
\toprule
\textbf{Variable} & \textbf{Pobre} & \textbf{No Pobre} & \textbf{Diferencia} \\
 & \textbf{(N=31,993)} & \textbf{(N=122,400)} & \\
 & Mean (SD)  & Mean (SD)  & Signif. \\
\midrule
\multicolumn{4}{@{}l}{\textit{Variables Continuas}} \\
\addlinespace[0.1cm]
Número personas hogar & 4.15 (2.03) & 3.09 (1.65) & *** \\
Cantidad cuartos & 3.04 (1.13) & 3.5 (1.24) & *** \\
Cantidad cuartos para dormir & 1.98 (0.84) & 1.99 (0.91) & * \\
Arriendo & 267,399.19 (2,288,990.88) & 494,500.74 (3,170,616.52) & *** \\
Edad promedio hogar & 30.6 (16.18) & 39.2 (16.63) & *** \\
Proporción mujeres hogar & 0.55 (0.24) & 0.52 (0.28) & *** \\
Régimen salud subsidio promedio hogar & 0.58 (0.34) & 0.32 (0.39) & *** \\
Informalidad promedio hogar & 0.03 (0.12) & 0.03 (0.15) & ** \\
Promedio de ocupados cuenta propia & 0.56 (0.46) & 0.37 (0.42) & *** \\
Ocupados promedio hogar & 0.31 (0.24) & 0.54 (0.32) & *** \\
Desocupados promedio hogar & 0.08 (0.17) & 0.05 (0.14) & *** \\
Jóvenes "ninis" promedio hogar & 0.09 (0.15) & 0.04 (0.12) & *** \\
Menores 5 años promedio hogar & 0.11 (0.16) & 0.05 (0.12) & *** \\
Menores 18 años promedio hogar & 0.37 (0.24) & 0.18 (0.21) & *** \\
Mayores 65 años promedio hogar & 0.1 (0.24) & 0.14 (0.28) & *** \\
Nivel educativo [ninguno] promedio hogar & 0.1 (0.2) & 0.04 (0.15) & *** \\
Nivel educativo [preescolar] promedio hogar & 0.03 (0.09) & 0.02 (0.07) & *** \\
Nivel educativo [primaria] promedio hogar & 0.33 (0.29) & 0.23 (0.3) & *** \\

\addlinespace[0.2cm]
\midrule
\multicolumn{4}{@{}l}{\textit{Variables Categóricas}} \\
\addlinespace[0.1cm]
\multicolumn{4}{@{}l}{\quad\textit{Tipo propiedad vivienda}} \\
\addlinespace[0.05cm]
\quad Propia, totalmente pagada & 28.0\% & 40.8\% & $-$*** \\
\quad Propia, en proceso de pago & 1.6\% & 3.7\% & $-$*** \\
\quad Arriendo & 42.7\% & 36.9\% & $+$*** \\
\quad Usufructo & 16.1\% & 15.3\% & $+$** \\
\quad Ocupación sin título legal & 11.4\% & 3.2\% & $+$*** \\
\quad Otro & 0.1\% & 0.1\% & $+$*** \\
\addlinespace[0.1cm]
\multicolumn{4}{@{}l}{\quad\textit{Urbano}} \\
\addlinespace[0.05cm]
\quad Rural & 14.5\% & 8.8\% & $+$*** \\
\quad Urbano & 85.5\% & 91.2\% & $-$*** \\
\addlinespace[0.2cm]
\midrule
\multicolumn{4}{@{}l}{\textit{Variables Dummy}} \\
\addlinespace[0.1cm]
Hogar jefe hogar ocupado & 13.3\% & 57.4\% & 44.1\% \\
Hogar jefe hogar desocupado & 2.0\% & 2.7\% & 0.6\% \\
Hogar jefe hogar mujer jefe & 9.7\% & 32.4\% & 22.6\% \\
Hogar jefe hogar regimén subsidiado & 15.4\% & 26.5\% & 11.1\% \\
Hogar jefe hogar cuenta propia & 9.9\% & 26.0\% & 16.1\% \\
\bottomrule
\end{tabularx}
\medskip
\tiny
\begin{minipage}{\textwidth}
\textit{Nota:} Las variables continuas se presentan mediante su media, desviación estándar (DE) y diferencia de medias mostrando su significancia. Las variables categóricas con múltiples niveles reportan proporciones por grupo (pobre) y se evalúa la significancia estadística mediante los residuos del test Chi-cuadrado, los signos refieren una mayor ($+$) o menor ($-$) probabilidad de ser pobre. Las variables dicotómicas (dummies) muestran la proporción en cada grupo sin prueba asociada. El total de observaciones es n = 154.393, con grupos definidos como 0 = No pobre y 1 = Pobre.)\\
Significancia estadística: *** $p < 0.001$, ** $p < 0.01$, * $p < 0.05$.
\end{minipage}
\end{table}